\section{Decoding results}
\label{sec:results}

\begin{table}[t]
\centering
\caption{Cross-subject decoding on ErrP-Bench: macro-participant AUROC (mean $\pm$ SD over training seeds;
$^{\ast}$: one seed). Bold: best per dataset. Rank: mean rank over datasets ($^{\dagger}$: not run on every
dataset). EA: per-participant Euclidean alignment. Pairwise tests of the best decoder against every other are
in Appendix~\ref{app:stats}.}
\label{tab:main}
\scriptsize
\setlength{\tabcolsep}{3pt}
\resizebox{\linewidth}{!}{\begin{tabular}{lllcccccc}
\toprule
Family & Model & EA & ErrP-HRI robot & ErrP-HRI cursor & BNCI 2015-013 & ERP CORE ERN & IntEr-HRI & Rank \\
\midrule
Riemannian & xDAWN+RG &  & 0.590\scriptsize{$^{\ast}$} & 0.823\scriptsize{$^{\ast}$} & 0.820\scriptsize{$^{\ast}$} & 0.945\scriptsize{$^{\ast}$} & 0.955\scriptsize{$^{\ast}$} & 15.6 \\
Riemannian & xDAWN+RG & \checkmark & 0.603\scriptsize{$^{\ast}$} & 0.851\scriptsize{$^{\ast}$} & 0.803\scriptsize{$^{\ast}$} & 0.950\scriptsize{$^{\ast}$} & 0.963\scriptsize{$^{\ast}$} & 13.4 \\
\midrule
compact & EEGNet &  & 0.570\scriptsize{$\pm$0.008} & 0.872\scriptsize{$\pm$0.002} & 0.823\scriptsize{$^{\ast}$} & 0.959\scriptsize{$\pm$0.002} & 0.974\scriptsize{$\pm$0.000} & 10.0 \\
compact & EEGNet & \checkmark & 0.606\scriptsize{$\pm$0.004} & 0.885\scriptsize{$\pm$0.006} & 0.803\scriptsize{$\pm$0.002} & 0.963\scriptsize{$\pm$0.000} & 0.947\scriptsize{$\pm$0.006} & 11.2 \\
compact & EEG-Inception &  & 0.591\scriptsize{$\pm$0.000} & 0.873\scriptsize{$\pm$0.011} & 0.849\scriptsize{$\pm$0.005} & 0.958\scriptsize{$\pm$0.002} & 0.970\scriptsize{$\pm$0.002} & 9.2 \\
compact & EEG-Inception & \checkmark & 0.632\scriptsize{$\pm$0.013} & 0.902\scriptsize{$\pm$0.002} & 0.809\scriptsize{$\pm$0.003} & 0.966\scriptsize{$\pm$0.001} & 0.967\scriptsize{$\pm$0.002} & 7.2 \\
compact & EEGNeX &  & 0.623\scriptsize{$\pm$0.000} & 0.878\scriptsize{$\pm$0.010} & 0.841\scriptsize{$\pm$0.001} & 0.963\scriptsize{$\pm$0.002} & 0.976\scriptsize{$\pm$0.003} & 5.8 \\
compact & EEGNeX & \checkmark & 0.622\scriptsize{$\pm$0.014} & 0.897\scriptsize{$\pm$0.005} & 0.809\scriptsize{$\pm$0.013} & 0.968\scriptsize{$\pm$0.000} & 0.971\scriptsize{$\pm$0.004} & 7.2 \\
compact & EEG-Conformer &  & 0.594\scriptsize{$\pm$0.004} & 0.792\scriptsize{$\pm$0.009} & 0.820\scriptsize{$\pm$0.003} & 0.950\scriptsize{$\pm$0.001} & 0.964\scriptsize{$\pm$0.004} & 14.0 \\
compact & EEG-Conformer & \checkmark & 0.597\scriptsize{$\pm$0.002} & 0.840\scriptsize{$\pm$0.001} & 0.788\scriptsize{$\pm$0.000} & 0.953\scriptsize{$\pm$0.002} & 0.960\scriptsize{$\pm$0.002} & 14.2 \\
compact & CTNet &  & 0.596\scriptsize{$\pm$0.016} & 0.873\scriptsize{$\pm$0.004} & 0.846\scriptsize{$\pm$0.023} & 0.959\scriptsize{$\pm$0.001} & 0.954\scriptsize{$\pm$0.000} & 10.8 \\
compact & CTNet & \checkmark & 0.617\scriptsize{$\pm$0.003} & 0.885\scriptsize{$\pm$0.005} & 0.820\scriptsize{$\pm$0.009} & 0.961\scriptsize{$\pm$0.002} & 0.944\scriptsize{$\pm$0.006} & 10.4 \\
\midrule
foundation & LaBraM &  & 0.554\scriptsize{$\pm$0.001} & 0.824\scriptsize{$\pm$0.002} & 0.801\scriptsize{$\pm$0.010} & 0.944\scriptsize{$\pm$0.001} & 0.955\scriptsize{$\pm$0.007} & 18.0 \\
foundation & LaBraM & \checkmark & 0.564\scriptsize{$\pm$0.002} & 0.851\scriptsize{$\pm$0.015} & 0.779\scriptsize{$\pm$0.014} & 0.952\scriptsize{$\pm$0.006} & 0.953\scriptsize{$\pm$0.007} & 16.8 \\
foundation & EEGPT &  & 0.582\scriptsize{$^{\ast}$} & 0.854\scriptsize{$^{\ast}$} & 0.812\scriptsize{$^{\ast}$} & -- & 0.957\scriptsize{$^{\ast}$} & 14.0$^{\dagger}$ \\
foundation & EEGPT & \checkmark & 0.623\scriptsize{$^{\ast}$} & 0.891\scriptsize{$^{\ast}$} & 0.778\scriptsize{$^{\ast}$} & -- & 0.920\scriptsize{$^{\ast}$} & 13.2$^{\dagger}$ \\
foundation & CBraMod &  & 0.554\scriptsize{$\pm$0.008} & 0.838\scriptsize{$\pm$0.004} & 0.819\scriptsize{$\pm$0.017} & 0.944\scriptsize{$\pm$0.002} & 0.959\scriptsize{$\pm$0.005} & 16.2 \\
foundation & CBraMod & \checkmark & 0.597\scriptsize{$\pm$0.001} & 0.845\scriptsize{$\pm$0.009} & 0.823\scriptsize{$\pm$0.002} & 0.952\scriptsize{$\pm$0.000} & 0.930\scriptsize{$\pm$0.006} & 13.8 \\
\midrule
ours & ErrP-Ens & both & 0.642\scriptsize{$\pm$0.003} & 0.906\scriptsize{$\pm$0.002} & \textbf{0.872\scriptsize{$\pm$0.004}} & 0.970\scriptsize{$\pm$0.001} & \textbf{0.981\scriptsize{$\pm$0.001}} & 1.6 \\
ours & ErrP-Ens-EA & \checkmark & \textbf{0.649\scriptsize{$\pm$0.001}} & \textbf{0.908\scriptsize{$\pm$0.003}} & 0.851\scriptsize{$\pm$0.002} & \textbf{0.971\scriptsize{$\pm$0.001}} & 0.974\scriptsize{$\pm$0.002} & 1.6 \\
ours & ErrP-Ens-EA (top-3) & \checkmark & 0.640\scriptsize{$\pm$0.008} & 0.905\scriptsize{$\pm$0.001} & 0.833\scriptsize{$\pm$0.001} & 0.969\scriptsize{$\pm$0.000} & 0.973\scriptsize{$\pm$0.001} & 4.0 \\
\bottomrule
\end{tabular}
}
\end{table}

Table~\ref{tab:main} reports every decoder on every dataset. Four results stand out.

\paragraph{The observed agent dominates decoder choice.} The same twelve participants watched both a robot and a
cursor commit the same kind of error. The best cross-subject decoder reaches \BestErrpHriCursor{} AUROC
(\BestModelErrpHriCursor{}) on cursor errors but only \BestErrpHriRobot{} (\BestModelErrpHriRobot{}) on robot
errors. Paired over participants, the cursor condition is higher by \CursorMinusRobot{} for \CursorGtRobot{}
participants (Wilcoxon $p=\CursorRobotP{}$). No decoder, pretrained or not, closes this gap: the
difference between the two conditions is larger than the spread of all nine decoders on either of them. A plausible
reason, which we do not test, is that a robot movement unfolds over time, so the moment at which it becomes an
error is less sharply defined than a cursor jump and the time-locked potential is smeared. For human--robot interaction, this is the
condition that matters, and it is the hardest one in the benchmark.

\paragraph{Pretraining does not buy cross-subject ErrP decoding.} The best foundation model minus the best
compact model averages \FMMinusCompactMean{} AUROC over datasets, and a foundation model is best on
\FMWinsDatasets{} datasets (per-dataset paired tests in Appendix~\ref{app:stats}). On ErrP-HRI robot,
LaBraM and CBraMod are among the weakest decoders. The decoders that do best are small and strongly
inductive: the Riemannian xDAWN pipeline, which needs seconds of CPU time, and EEGNet-style networks. One
interpretation is that ErrPs are low-dimensional, time-locked waveforms: a model that already encodes spatial
filtering and a short temporal template has little left to gain from generic pretraining, whereas a large
backbone fine-tuned on a few thousand trials from about ten participants can fit the training participants
too closely.

\paragraph{Unsupervised alignment helps when the paradigm is consistent.} EA changes AUROC by \EAMeanDelta{} on
average over all model--dataset pairs, improving \EAPosFrac{} of them (Table~\ref{tab:ea}). Its effect depends
on the dataset: \EADeltaErrpHriCursor{} on ErrP-HRI cursor (\EASigPosErrpHriCursor{} of \EANModelsErrpHriCursor{}
models significantly better) but \EADeltaInterHri{} on IntEr-HRI (\EASigNegInterHri{} significantly worse).
EA removes between-participant differences in spatial covariance; our hypothesis for the negative cases is that
when that covariance is shaped by the task itself (in IntEr-HRI, the arm movement and the ball squeeze),
whitening it also removes signal. The averages over families are \EAScratchDelta{} for networks trained from scratch and \EAFMDelta{}
for foundation models.

\paragraph{Label-free ensembling.} ErrP-Ens averages, per participant, the within-participant ranks of the
aligned decoders. It needs no labels from the new participant and no training. Its mean rank over datasets is
\EnsAllMeanRank{} (all members) and \EnsTopMeanRank{} (validation top-3), among up to \NVariantsMax{} variants.
Against the best single decoder of each dataset (chosen with hindsight on the test participants, so the
comparison favours the single model), ErrP-Ens (all) differs by \EnsAllVsBestMean{} on average, is higher on
\EnsAllBeatsBest{} datasets, and is significantly better on \EnsAllSigBetter{} and significantly worse on
\EnsAllSigWorse{} after Holm correction.

\begin{table}[t]
\centering
\caption{Effect of Euclidean alignment: AUROC with EA minus without, averaged over seeds and participants.
$^{\ast}$: paired Wilcoxon $p<0.05$ over participants (uncorrected).}
\label{tab:ea}
\scriptsize
\resizebox{\linewidth}{!}{\begin{tabular}{lcccccc}
\toprule
Model & ErrP-HRI robot & ErrP-HRI cursor & BNCI 2015-013 & ERP CORE ERN & IntEr-HRI & Mean $\Delta$ \\
\midrule
xDAWN+RG & +0.013 & +0.028$^{\ast}$ & -0.017 & +0.006 & +0.008 & +0.007 \\
EEGNet & +0.036$^{\ast}$ & +0.013 & -0.020 & +0.005$^{\ast}$ & -0.027$^{\ast}$ & +0.002 \\
EEG-Inception & +0.041 & +0.028$^{\ast}$ & -0.041$^{\ast}$ & +0.009$^{\ast}$ & -0.003 & +0.007 \\
EEGNeX & -0.001 & +0.018 & -0.033 & +0.006$^{\ast}$ & -0.005 & -0.003 \\
EEG-Conformer & +0.002 & +0.048$^{\ast}$ & -0.032 & +0.002 & -0.004 & +0.004 \\
CTNet & +0.021$^{\ast}$ & +0.012 & -0.026 & +0.003 & -0.010 & +0.000 \\
LaBraM & +0.009 & +0.026$^{\ast}$ & -0.021 & +0.008 & -0.002 & +0.004 \\
EEGPT & +0.041$^{\ast}$ & +0.037$^{\ast}$ & -0.034 & -- & -0.037$^{\ast}$ & +0.002 \\
CBraMod & +0.043$^{\ast}$ & +0.007 & +0.004 & +0.007 & -0.029$^{\ast}$ & +0.006 \\
\bottomrule
\end{tabular}
}
\end{table}

\paragraph{IntEr-HRI measures more than an ErrP.} Decoders reach AUROC near \BestInterHri{} on IntEr-HRI. In
about 90\% of its error epochs the participant's ball squeeze, a motor response to the error, falls inside the
window. Restricted to the first 0.5\,s, before almost all squeezes (0.8\% occur earlier), EEGNet still reaches
\EarlyEEGNetNone{} and xDAWN \EarlyXdawnNone{}. A large part of the discriminative signal therefore precedes
the squeeze, but the full-window scores should be read as an upper bound for an error-plus-response signal.
